% BHU PYQ -- Professional Exam Template
\documentclass[11pt,a4paper]{article}

\usepackage[a4paper,top=2.5cm,bottom=2.5cm,left=2.8cm,right=2.8cm,headheight=14pt]{geometry}
\usepackage{amsmath,amssymb}
\usepackage[T1]{fontenc}
\usepackage[utf8]{inputenc}
\usepackage{lmodern}
\usepackage{microtype}
\usepackage{fancyhdr}
\usepackage{enumitem}
\usepackage{hyperref}
\usepackage{lastpage}
\usepackage{parskip}

\hypersetup{colorlinks=true,urlcolor=black,linkcolor=black,
  pdftitle={B.Sc. (Hons.) Botany Sem-VI BOB-601 -- BHU PYQ}}

\pagestyle{fancy}
\fancyhf{}
\renewcommand{\headrulewidth}{0.4pt}
\renewcommand{\footrulewidth}{0.4pt}
\fancyhead[L]{\small   \textit{Banaras Hindu University}}
\fancyhead[R]{\small   \textit{Botany Paper: BOB-601}}
\fancyfoot[C]{\small Page  \thepage\ of \pageref{LastPage}}


\newlist{parts}{enumerate}{2}
\setlist[parts,1]{label=(\alph*),leftmargin=*,itemsep=4pt,topsep=3pt}
\setlist[parts,2]{label=(\roman*),leftmargin=*,itemsep=2pt,topsep=2pt}

\newcommand{\pts}[1]{\hfill   \textit{[#1]}}


\begin{document}


\begin{center}
  {\large\bfseries BANARAS HINDU UNIVERSITY}\\[2pt]
  {
\normalsize B.Sc. (Hons.) Semester VI Examination, 2013-14}\\[6pt]
  \rule{0.8\linewidth}{0.5pt}\\[6pt]
  {\large\bfseries Botany}\\[4pt]
  {\large\bfseries Paper: BOB-601 --- Plant Metabolism, Biochemistry and Biotechnology}\\[4pt]
  {
\normalsize Time: Three hours \hfill Full Marks: 70}
\end{center}

\rule{\linewidth}{0.4pt}

\smallskip



\noindent   \textit{Note: Attempt five questions in all, selecting two questions each from Section-A and B. Question No. 9 is compulsory. The figures in the right-hand margin indicate marks.}

\bigskip


\begin{center}
   \textbf{SECTION -- A}
\end{center}

\medskip


\noindent   \textbf{Question 1.} What are co-enzymes? How do the allosteric enzymes differ from the enzymes following Michaelis-Menten kinetics?\pts{15}

\medskip


\noindent   \textbf{Question 2.} Describe the pathway of sucrose biosynthesis in plants.\pts{15}

\medskip


\noindent   \textbf{Question 3.} Write critical notes on any two of the following:\pts{$7.5  \times 2 = 15$}

\begin{parts}
  \item Ammonia assimilation
  \item Mechanism of $ \text{O}_2$ regulation during $ \text{N}_2$ fixation
  \item Biosynthesis of cellulose
\end{parts}

\medskip


\noindent   \textbf{Question 4.} Write short notes on any two of the following:\pts{$7.5  \times 2 = 15$}

\begin{parts}
  \item Nitrogenase enzyme
  \item Activation and assimilation of phosphorus
  \item Isoenzymes
\end{parts}

\bigskip

\begin{center}
   \textbf{SECTION -- B}
\end{center}

\medskip


\noindent   \textbf{Question 5.} What are auxins? Describe their biosynthetic pathways and mode of action.\pts{15}

\medskip


\noindent   \textbf{Question 6.} What is DNA replication? Describe the process of DNA replication and its control mechanisms.\pts{15}

\medskip


\noindent   \textbf{Question 7.} Write short notes on any two of the following:\pts{$7.5  \times 2 = 15$}

\begin{parts}
  \item Cloning vectors
  \item B-DNA
  \item Recombinant DNA technology
\end{parts}

\medskip


\noindent   \textbf{Question 8.} Briefly describe any two of the following:\pts{$7.5  \times 2 = 15$}

\begin{parts}
  \item Transcription
  \item Abscisic acid
  \item Post-translational modification of proteins
\end{parts}

\bigskip

\begin{center}
   \textbf{SECTION -- C}
\end{center}

\medskip


\noindent   \textbf{Question 9.} Fill in the blanks or answer the following in not more than 50 words each:\pts{$1  \times 10 = 10$}

\begin{parts}
  \item What are energy-rich phosphorus compounds?
  \item Cell division can be stimulated by \makebox[3cm]{\hrulefill} hormone.
  \item Define prosthetic groups.
  \item \makebox[3cm]{\hrulefill} is the site of lipid synthesis in plants.
  \item EF-Ts is an \makebox[3cm]{\hrulefill} factor.
  \item What is the role of topoisomerase?
  \item \makebox[3cm]{\hrulefill} gene encodes DNA polymerase I.
  \item The number of base pair per turn of A-DNA helix is \makebox[2cm]{\hrulefill}.
  \item What is a shuttle vector?
  \item \makebox[3cm]{\hrulefill} is an initiation codon.
\end{parts}

\vfill
\rule{\linewidth}{0.4pt}

\begin{center}\small
\href{https://bhu-exam.vercel.app/}{Click Here}\\[4pt]\href{https://me-aryan.vercel.app/}{Click Here}\\[4pt]\href{https://quashan-me.vercel.app/}{Click Here}
\end{center}

\end{document}
